\documentclass[openany]{book}

\usepackage[margin=1in]{geometry}
\usepackage{amsmath,amsfonts,amsthm, amssymb}
\usepackage{yhmath}
\usepackage{mathrsfs}
\usepackage{mathtools}
\usepackage{xcolor}
\usepackage{graphicx}
\usepackage{comment}
\usepackage{tikz-cd}
\usepackage{quiver}
\renewcommand{\familydefault}{ppl}
\newcommand{\tr}{\text{tr}}
\newcommand{\R}{\mathbb{R}}
\newcommand{\E}{\mathbb{E}}
\newcommand{\Z}{\mathbb{Z}}
\newcommand{\C}{\mathbb{C}}
\newcommand{\F}{\mathbb{F}}
\newcommand{\la}{\langle}
\newcommand{\ra}{\rangle}
\newcommand{\colim}{\text{colim}}
\DeclareMathOperator{\im}{im}
\let\oldemptyset\emptyset
\let\emptyset\varnothing
\newcommand{\tor}{\text{Tor}}
\newcommand{\id}{\text{id}}
\newcommand{\ext}{\text{Ext}}
\newcommand{\ptop}{\text{PTop}}
\newcommand{\pt}{\text{pt}}
\newcommand{\ach}{\text{Ach}}
\newcommand{\Q}{\mathbb{Q}}
\newcommand{\gal}{\text{Gal}}


\usepackage{thmtools,thm-restate}

% Fixing mdframed skip below
% See https://tex.stackexchange.com/a/292090/143086
\usepackage[framemethod=TikZ]{mdframed}
\usepackage{xpatch}
\makeatletter
\xpatchcmd{\endmdframed}
	{\aftergroup\endmdf@trivlist\color@endgroup}
	{\endmdf@trivlist\color@endgroup\@doendpe}
	{}{}
\makeatother

\definecolor{huilightpink}{HTML}{fff2fe}
\definecolor{huidarkpink}{HTML}{d955b7}
\declaretheoremstyle[
	mdframed={
		backgroundcolor=huilightpink,
		linecolor=huidarkpink,
		rightline=false,
		topline=false,
		bottomline=false,
		linewidth=2pt,
		innertopmargin=5pt,
		innerbottommargin=8pt,
		innerleftmargin=8pt,
		leftmargin=-2pt,
		skipbelow=2pt,
		nobreak
	},
	headfont=\normalfont\bfseries\color{huidarkpink}
]{huipinkbox}
\declaretheorem[style=huipinkbox,name=Theorem,within=chapter]{thm}
\declaretheorem[style=huipinkbox,name=Theorem,sibling=thm]{theorem}




\begin{comment}
\definecolor{huilightyellow}{HTML}{fff5d6}
\definecolor{huidarkyellow}{HTML}{fcad03}
\declaretheoremstyle[
	mdframed={
		backgroundcolor=huilightyellow,
		linecolor=huidarkyellow,
		rightline=false,
		topline=false,
		bottomline=false,
		linewidth=2pt,
		innertopmargin=5pt,
		innerbottommargin=8pt,
		innerleftmargin=8pt,
		leftmargin=-2pt,
		skipbelow=2pt,
		nobreak
	},
	headfont=\normalfont\bfseries\color{huidarkyellow}
]{huiyellowbox}
\declaretheorem[style=huiyellowbox,name=Proposition,within=chapter]{prop}
\end{comment}



\definecolor{huilightpurple}{HTML}{faf2ff}
\definecolor{huidarkpurple}{HTML}{912ed9}
\declaretheoremstyle[
	mdframed={
		backgroundcolor=huilightpurple,
		linecolor=huidarkpurple,
		rightline=false,
		topline=false,
		bottomline=false,
		linewidth=2pt,
		innertopmargin=5pt,
		innerbottommargin=8pt,
		innerleftmargin=8pt,
		leftmargin=-2pt,
		skipbelow=2pt,
		nobreak
	},
	headfont=\normalfont\bfseries\color{huidarkpurple}
]{huipurplebox}
\declaretheorem[style=huipurplebox,name=Proposition,within=chapter]{prop}



% \definecolor{huilightpurple}{HTML}{faf2ff}
% \definecolor{huidarkpurple}{HTML}{912ed9}
% \declaretheoremstyle[
% 	mdframed={
% 		backgroundcolor=huilightpurple,
% 		linecolor=huidarkpurple,
% 		rightline=false,
% 		topline=false,
% 		bottomline=false,
% 		linewidth=2pt,
% 		innertopmargin=5pt,
% 		innerbottommargin=8pt,
% 		innerleftmargin=8pt,
% 		leftmargin=-2pt,
% 		skipbelow=2pt,
% 		nobreak
% 	},
% 	headfont=\normalfont\bfseries\color{huidarkpurple}
% ]{huipurplebox}
\declaretheorem[style=huipurplebox,name=Lemma,within=chapter]{lem}


\definecolor{lightpink}{HTML}{f0f6fc}
\definecolor{darkpink}{HTML}{2c72b8}
\declaretheoremstyle[
	mdframed={
		backgroundcolor=lightpink,
		linecolor=darkpink,
		rightline=false,
		topline=false,
		bottomline=false,
		linewidth=2pt,
		innertopmargin=5pt,
		innerbottommargin=8pt,
		innerleftmargin=8pt,
		leftmargin=-2pt,
		skipbelow=2pt,
		nobreak
	},
	headfont=\normalfont\bfseries\color{darkpink}
]{pinkbox}
\declaretheorem[style=pinkbox,name=Definition,within=chapter]{defn}


\definecolor{huilightblue}{HTML}{edf9ff}
\definecolor{huidarkblue}{HTML}{4b79db}
\declaretheoremstyle[
	mdframed={
		backgroundcolor=huilightblue,
		linecolor=huidarkblue,
		rightline=false,
		topline=false,
		bottomline=false,
		linewidth=2pt,
		innertopmargin=5pt,
		innerbottommargin=8pt,
		innerleftmargin=8pt,
		leftmargin=-2pt,
		skipbelow=2pt,
		nobreak
	},
	headfont=\normalfont\bfseries\color{huidarkblue}
]{huiblueblox}
\declaretheorem[style=huiblueblox,name=Example,within=chapter]{example}



% \definecolor{huilightblue}{HTML}{edf9ff}
% \definecolor{huidarkblue}{HTML}{4b79db}
% \declaretheoremstyle[
% 	mdframed={
% 		backgroundcolor=huilightblue,
% 		linecolor=huidarkblue,
% 		rightline=false,
% 		topline=false,
% 		bottomline=false,
% 		linewidth=2pt,
% 		innertopmargin=5pt,
% 		innerbottommargin=8pt,
% 		innerleftmargin=8pt,
% 		leftmargin=-2pt,
% 		skipbelow=2pt,
% 		nobreak
% 	},
% 	headfont=\normalfont\bfseries\color{huidarkblue}
% ]{huiblueblox}
% \declaretheorem[style=huiblueblox,name=Example,within=chapter]{example}

% \declaretheoremstyle[
% 	mdframed={
% 		backgroundcolor=huilightblue,
% 		linecolor=huidarkblue,
% 		rightline=false,
% 		topline=false,
% 		bottomline=false,
% 		linewidth=2pt,
% 		innertopmargin=5pt,
% 		innerbottommargin=8pt,
% 		innerleftmargin=8pt,
% 		leftmargin=-2pt,
% 		skipbelow=2pt,
% 		nobreak
% 	},
% 	headfont=\normalfont\bfseries\color{huidarkblue}
% ]{huiblueblox}
\declaretheorem[style=huiblueblox,name=Problem,within=chapter]{prob}



% \declaretheoremstyle[
% 	mdframed={
% 		backgroundcolor=huilightblue,
% 		linecolor=huidarkblue,
% 		rightline=false,
% 		topline=false,
% 		bottomline=false,
% 		linewidth=2pt,
% 		innertopmargin=5pt,
% 		innerbottommargin=8pt,
% 		innerleftmargin=8pt,
% 		leftmargin=-2pt,
% 		skipbelow=2pt,
% 		nobreak
% 	},
% 	headfont=\normalfont\bfseries\color{huidarkblue}
% ]{huiblueblox}
\declaretheorem[style=huiblueblox,name=Exercise,within=chapter]{exer}
\declaretheorem[style=huipinkbox, name=Corollary, within=chapter]{cor}





















\newcommand{\nirwarnsymbol}{%
	\begin{tikzpicture}[baseline=(x.base)]
		\draw[rounded corners=.01em] (-.05em,-1.07em)rectangle(.05em,.78em);
		\draw[fill=white,rounded corners=1.3] (0,.75em)--(.75em,0)--(0,-.75em)--(-.75em,0)--cycle;
		\draw[line width=0.2mm, line cap=round](-.4em,-1.07em)--(.4em,-1.07em);
		\node(x) at (0,0em) {};
		% Thank you https://tex.stackexchange.com/a/262510
		\draw[
			line cap=but,
			line join=round,
			x=.5em,
			line width=0.5mm,
			y=1*(height("Z")-\pgflinewidth)*(1-sin(10)),
			rotate=-10,
			rounded corners=1.5pt,
		](-0.57, 0.57) -- (0.57, 0.57) -- (-0.57, -0.57) -- (0.57, -0.57);
	\end{tikzpicture}%
}

%%%%%%%%%%%%%%%%%%%%%%%%%%%%%%%%%%%%%%%%%%%% MARGINS
\usepackage{marginnote}
% Thank you https://tex.stackexchange.com/a/472882
% Makes marginnotes always appear on the left, apparently
%
\makeatletter
\long\def\@mn@@@marginnote[#1]#2[#3]{%
	\begingroup
		\ifmmode\mn@strut\let\@tempa\mn@vadjust\else
			\if@inlabel\leavevmode\fi
			\ifhmode\mn@strut\let\@tempa\mn@vadjust\else\let\@tempa\mn@vlap\fi
		\fi
		\@tempa{%
			\vbox to\z@{%
				\vss
				\@mn@margintest
				\if@reversemargin\if@tempswa
						\@tempswafalse
					\else
						\@tempswatrue
				\fi\fi

					\llap{%
						\vbox to\z@{\kern\marginnotevadjust\kern #3
							\vbox to\z@{%
								\hsize\marginparwidth
								\linewidth\hsize
								\kern-\parskip
								%\mn@parboxrestore
								\marginfont\raggedleftmarginnote\strut\hspace{\z@}%
								\ignorespaces#1\endgraf
								\vss
							}%
							\vss
						}%
						\if@mn@verbose
							\PackageInfo{marginnote}{xpos seems to be \@mn@currxpos}%
						\fi
						\begingroup
							\ifx\@mn@currxpos\relax\else\ifx\@mn@currpos\@empty\else
									\kern\@mn@currxpos
							\fi\fi
							\ifx\@mn@currpage\relax
								\let\@mn@currpage\@ne
							\fi
							\if@twoside\ifodd\@mn@currpage\relax
									\kern-\oddsidemargin
								\else
									\kern-\evensidemargin
								\fi
							\else
								\kern-\oddsidemargin
							\fi
							\kern-1in
						\endgroup
						\kern\marginparsep
					}%
			}%
		}%
	\endgroup
}
\makeatother
%
% Mostly for todonotes
\renewcommand{\marginpar}{\marginnote}
%%%%%%%%%%%%%%%%%%%%%%%%%%%%%%%%%%%%%%%%%%%% /MARGINS

\definecolor{nirlightred}{RGB}{250, 220, 220}
\definecolor{nirdarkred}{HTML}{f40000}
\declaretheoremstyle[
	mdframed={
		backgroundcolor=nirlightred,
		linecolor=nirdarkred,
		rightline=false,
		topline=false,
		bottomline=false,
		linewidth=2pt,
		innertopmargin=5pt,
		innerbottommargin=8pt,
		innerleftmargin=8pt,
		leftmargin=-2pt,
		skipbelow=2pt,
		nobreak
	},
	headfont=\normalfont\bfseries\color{nirdarkred}
]{nirredbox}

% \makeatletter
% \declaretheorem[
% 	style=nirredbox,
% 	name=Warning,
% 	sibling=thm,
% 	% without \leavevmode, the first item in a list gets misformatted
% 	postheadhook={\leavevmode\marginnote{\nirwarnsymbol}[-3pt]%
% 	\ifthmt@thisistheone% restatable makes alignment weird
% 		\hspace{-2.2pt}%
% 	\fi}
% ]{warn}
% \makeatother

\newcommand{\nirideasymbol}{%
	\begin{tikzpicture}[baseline=(x.base)]
		\draw[rounded corners=.01em] (-.05em,-1.07em)rectangle(.05em,.78em);
		\draw[fill=white,rounded corners=1.3] (0,.75em)--(.75em,0)--(0,-.75em)--(-.75em,0)--cycle;
		\draw[line width=0.2mm, line cap=round](-.4em,-1.07em)--(.4em,-1.07em);
		\node(x) at (0,0em) {};
		\node at (0,0em) {{\textbf{!}}};
	\end{tikzpicture}%
}
\renewcommand{\nirwarnsymbol}{%
	\begin{tikzpicture}[baseline=(x.base)]
		\draw[rounded corners=.01em] (-.05em,-1.07em)rectangle(.05em,.78em);
		\draw[fill=white,rounded corners=1.3] (0,.75em)--(.75em,0)--(0,-.75em)--(-.75em,0)--cycle;
		\draw[line width=0.2mm, line cap=round](-.4em,-1.07em)--(.4em,-1.07em);
		\node(x) at (0,0em) {};
		% Thank you https://tex.stackexchange.com/a/262510
		\draw[
			line cap=but,
			line join=round,
			x=.5em,
			line width=0.5mm,
			y=1*(height("Z")-\pgflinewidth)*(1-sin(10)),
			rotate=-10,
			rounded corners=1.5pt,
		](-0.57, 0.57) -- (0.57, 0.57) -- (-0.57, -0.57) -- (0.57, -0.57);
	\end{tikzpicture}%
}
\makeatletter
\declaretheorem[
	style=nirredbox,
	name=Idea,
	sibling=thm,
	% without \leavevmode, the first item in a list gets misformatted
	postheadhook={\leavevmode\marginnote{\nirideasymbol}[-3pt]%
	\ifthmt@thisistheone% restatable makes alignment weird
		\hspace{-2.2pt}%
	\fi}
]{idea}

\declaretheorem[
	style=nirredbox,
	name=Warning,
	sibling=thm,
	% without \leavevmode, the first item in a list gets misformatted
	postheadhook={\leavevmode\marginnote{\nirwarnsymbol}[-3pt]%
	\ifthmt@thisistheone% restatable makes alignment weird
		\hspace{-2.2pt}%
	\fi}
]{warn}
\makeatother

\title{Calc III Sections
\\ 
\vspace{0.4cm}
\large Fall 2025}




\date{\today}
\author{Hui Sun}


\begin{document}

\maketitle

% \tableofcontents
\newpage


\begin{center}
    \Large Calc III-Week 7 (10/6-10/10)
\end{center}

\renewcommand\thesection{\arabic{section}}

\noindent
Topics: (1) quadratic form, (2) constraint. 




% \begin{defn}[critical point]
%     Let $f:U\subset\R^n\to\R$, a point $x_0\in U$ is a \textbf{critical point} of $f$ if either $f$ is not differentiable at $x_0$, or $Df(x_0)=0$. A critical point that is not a local extremum is called a saddle point.
% \end{defn}

\begin{defn}[quadratic function]
    A function $g:\R^n\to\R$ is called a \textbf{quadratic function} if it is given by 
    \begin{equation*}
        g(h_1,\dots, h_n)=\sum_{i,j=1}^na_{ij}h_ih_j
    \end{equation*}
    where $(a_{ij})$ is an $n\times n$ matrix. We can also write $g$ as follows:
    \begin{equation*}
        g(h_1,\dots,h_n)=[h_1,\dots,h_n]\begin{bmatrix}
            a_{11}&\dots &a_{1n}\\
            \vdots&\ddots &\vdots\\
            a_{n_1}&\dots&a_{nn}
        \end{bmatrix}\begin{bmatrix}
            h_1\\
            \vdots\\
            h_n
        \end{bmatrix}
    \end{equation*}
    An example of a quadratic function is defined by 
    \begin{equation*}
        Hf(x)(h)=\frac{1}{2}\begin{bmatrix}
            h_1&\dots&h_n
        \end{bmatrix}Hf(x)\begin{bmatrix}
            h_1\\
            \dots\\
            h_n
        \end{bmatrix}
    \end{equation*}
    where $h=(h_1,\dots, h_n)$, and $Hf(x)$ is the Hessian matrix of $f:\R^n\to\R^n$.
\end{defn}

% \begin{defn}[Hessian matrix]
%     Let $f:U\subset\R^n\to\R$, and suppose all the second-order partial derivatives $\frac{\partial^2 f}{\partial x_i\partial x_j}$ exist, then the Hessian matrix of $f$ is the $n\times n$ matrix given by 
%     \begin{equation*}
%         Hf=\begin{bmatrix}
%             \frac{\partial^2 f}{\partial x_1\partial x_1}&\dots&\frac{\partial^2f}{\partial x_1\partial x_n}\\
%             \vdots&\ddots&\vdots\\
%             \frac{\partial^2f}{\partial x_nx_1}&\dots&\frac{\partial^2f}{\partial x_n\partial x_n}
%         \end{bmatrix}
%     \end{equation*}
%     The Hessian as a quadratic function is defined by 
%     \begin{equation*}
%         Hf(x)(h)=\frac{1}{2}\begin{bmatrix}
%             h_1&\dots&h_n
%         \end{bmatrix}Hf(x)\begin{bmatrix}
%             h_1\\
%             \dots\\
%             h_n
%         \end{bmatrix}
%     \end{equation*}
%     where $h=(h_1,\dots,h_n)$.
% \end{defn}

% \begin{defn}[degenerate/nondegenerate points]
%     Let $f:U\subset\R^2\to\R$ be of $C^2$, let $(x_0,y_0)$ be a critical point. We define the \textbf{discriminant}, \boldmath $\mathcal{D}$,\unboldmath of the Hessian by
%     \begin{equation*}
%         \mathcal{D}=\det (Hf)=\left(\frac{\partial^2f}{\partial x^2}\right)\left(\frac{\partial^2f}{\partial y^2}\right)-\left(\frac{\partial^2f}{\partial x\partial y}\right)^2
%     \end{equation*}
%     If $\mathcal{D}\neq 0$, the critical point $(x_0,y_0)$ is called \textbf{nondegenerate}; if $\mathcal{D}=0$, the point $(x_0,y_0)$ is called \textbf{degenerate}.
% \end{defn}

\begin{defn}[positive, negative-definite]
    A quadratic function $g:\R^n\to\R$ is called \textbf{positive-definite} if $g(h)\geq 0$ for all $h\in\R^n$ and $g(h)=0$ implies $h=0$. Similarly, $g$ is \textbf{negative-definite} if $g(h)\leq 0$ for all $h\in\R^n$ and $g(h)=0$ implies $h=0$. (The matrix is positive-definite iff it is symmetric $A^T=A$ and the eigenvalues are nonnegative).
\end{defn}


% \begin{defn}[global extremum]
%     Let $f:A\to\R$ be a function defined on $A\subset\R^2$ or $A\subset\R^3$. A point $x_0\in A$ is said to be an \textbf{absolute maximum} if $f(x_0)\geq f(x)$ for all $x\in A$. Similarly, $x_0$ is an \textbf{absolute minimum} if $f(x_0)\leq f(x)$ for all $x\in A$. 
% \end{defn}

\begin{defn}[bounded set]
    A set $A\subset\R^n$ is said to be \textbf{bounded} if there is a number $M>0$ such that $\|x\|\leq M$ for all $x\in A$. 
\end{defn}




% \begin{prop}[extremums are critical points]
%     Let $f:U\subset\R^n\to\R$ be differentiable, where $U$ is open. If $x_0$ is a local extremum, then $Df(x_0)=0$. 
% \end{prop}


% \begin{prop}[extremum]
%     Let $f\colon U\subset\R^n\to\R$ be in $C^3$, and $x_0$ is a critical point of $f$. If the Hessian $Hf(x_0)$ is positive-definite, then $x_0$ is a local minimum of $f$; if $Hf(x_0)$ is negative-definite, then $x_0$ is a local maximum.
% \end{prop}




% \begin{prop}[local minimum]
%     Let $f(x,y)$ be of $C^2$, and $U$ is open in $\R^2$. A point $(x_0,y_0)$ is a strict local \textbf{minimum} of $f$ if the following conditions hold:
%     \begin{enumerate}
%         \item \begin{equation*}
%             \frac{\partial f}{\partial x}(x_0,y_0)=\frac{\partial f}{\partial y}(x_0,y_0)=0
%         \end{equation*}
%         \item \begin{equation*}
%             \mathcal{D}(x_0,y_0)>0
%         \end{equation*}
%         where $\mathcal{D}$ is the \textbf{discriminant} of the Hessian, defined by 
%         \begin{equation*}
%             \mathcal{D}=\det (Hf)=\left(\frac{\partial^2 f}{\partial x^2}\right)\left(\frac{\partial^2 f}{\partial y^2}\right)-\left(\frac{\partial^2 f}{\partial x\partial y}\right)^2
%         \end{equation*}
%         where $Hf$ is the $2\times 2$ Hessian matrix.
%         \item  \begin{equation*}
%             \frac{\partial^2f}{\partial x^2}(x_0,y_0)>0
%         \end{equation*}
%     \end{enumerate}
%     If $ \frac{\partial^2f}{\partial x^2}(x_0,y_0)<0$ in 3, then it becomes a local maximum.
% \end{prop}


% \begin{prop}[local maximum]
%     Let $f(x,y)$ be of $C^2$, and $U$ is open in $\R^2$. A point $(x_0,y_0)$ is a strict local \textbf{maximum} of $f$ if the following conditions hold:
%     \begin{enumerate}
%         \item \begin{equation*}
%             \frac{\partial f}{\partial x}(x_0,y_0)=\frac{\partial f}{\partial y}(x_0,y_0)=0
%         \end{equation*}
%         \item \begin{equation*}
%             \mathcal{D}(x_0,y_0)>0
%         \end{equation*}
%         where $\mathcal{D}$ is the {discriminant} of the Hessian, defined above.
%         \item  \begin{equation*}
%             \frac{\partial^2f}{\partial x^2}(x_0,y_0)<0
%         \end{equation*}
%     \end{enumerate}
% \end{prop}

% \begin{prop}[saddle points]
%     Let $f(x,y):U\subset\R^2\to \R$ be of  $C^2$, if $\frac{\partial f}{\partial x}(x_0,y_0)=\frac{\partial f}{\partial y}(x_0,y_0)=0$, and $\mathcal{D}(x_0,y_0)<0$, where $\mathcal{D}$ is the discriminant, then the critical point $(x_0,y_0)$ is a saddle point, i.e., neither a maximum or a minimum.
% \end{prop}






\begin{prop}[continuous functions attain extremum on closed bounded sets]
    Let $f:D\to\R$ be continuous, where $D$ is closed and bounded in $\R^n$. Then $f$ assumes its absolute maximum and absolute minimum values at some point $x_0, x_1\in D$.
\end{prop}



\begin{prop}
        If $f$ has a maximum or minimum
        at $x_0$ when constrained to a surface $S$, then $\nabla f(x)$ is perpendicular to $S$ at $x_0$.
\end{prop}
Consequence of Proposition 2:
\begin{prop}
    Let \( f, g : U \subset \mathbb{R}^n \to \mathbb{R} \) and \(\vec{x_0} \in U\) such that \( g({x_0}) = c \). Let \( \mathcal{L}_c \) be the level set for \( g \) with value \( c \) and assume  
\[
\nabla g(\vec{x_0}) \neq \vec{0}.
\] 
If \( f \) restricted to \( x_c \) has a local minimum or maximum on \( \mathcal{L}_c \) at \( {x_0} \), then there exists \( \lambda \in \mathbb{R} \) such that  
\[
\nabla f({x_0}) = \lambda \nabla g({x_0}).
\]
\end{prop}

\begin{prob}
    Is the following matrix positive-definite?
    \begin{equation*}
        \begin{bmatrix}
            1&0\\
            0&-1
        \end{bmatrix}
    \end{equation*}
\end{prob}
\begin{proof}
    It is not! Consider the vector $(0,1)$, we have 
    \begin{equation*}
        [0, 1]
\begin{bmatrix}
        1&0\\
        0&-1
    \end{bmatrix}\begin{bmatrix}
        0\\
        1
    \end{bmatrix}
    =[0,1]\begin{bmatrix}
        0\\
        -1
    \end{bmatrix}=-1
\end{equation*}
\end{proof}


\begin{prob}
    Find the critical point of $f(x,y)=y+x\sin y$ and classify whether it is a local max/min or a saddle point.
\end{prob}
\begin{proof}
    We compute 
    \begin{equation*}
        f_x=\sin y, \quad f_y=1+x\cos y
    \end{equation*}
    Setting them both to $0$ gives 
    \begin{equation*}
        (1, n\pi), \text{ when $n$ is odd}, \quad (-1,n\pi) \text{ when $n$ is even}
    \end{equation*}
    Now we compute the discriminant:
    \begin{equation*}
        f_{xx}=0, \quad f_{xy}=f_{yx}=\cos y, \quad f_{yy}=-x\sin y
    \end{equation*}
    Thus 
    \begin{equation*}
        \mathcal{D}=\det\begin{bmatrix}
            0&\cos y\\
            \cos y&-x\sin y
        \end{bmatrix}=-\cos^2y
    \end{equation*}
    Thus $\mathcal{D}<0$ for all the critical points, hence they are saddle points!
\end{proof}


\begin{prob}[Marsden-Tromba, III. 2]
    Let $f(x,y,z)=x-y+z$, find the extremum of $f$ subject to the constraint $x^2+y^2+z^2=2$.
\end{prob}
\begin{proof}
    We compute the gradient of $f$ and $g(x,y,z)=x^2+y^2+z^2$:
    \begin{equation*}
        \nabla f(x,y,z)=\begin{bmatrix}
            1&-1&1
        \end{bmatrix}\quad \nabla g(x,y,z)=\begin{bmatrix}
            2x&2y&2z
        \end{bmatrix}
    \end{equation*}
    By the proposition above, we need to have 
    \begin{align*}
        \lambda&=2x\\
        -\lambda&=2y\\
        \lambda&=2z
    \end{align*}
    Equating all $\lambda$ gives 
    \begin{equation*}
        x=z=-y
    \end{equation*}
    Plugging in the constraint we get 
    \begin{equation*}
        3x^2=2
    \end{equation*}
    Hence two critical points 
    \begin{equation*}
        \left(\sqrt{\frac{2}{3}}, -\sqrt{\frac{2}{3}}, \sqrt{\frac{2}{3}}\right), \quad \left(-\sqrt{\frac{2}{3}}, \sqrt{\frac{2}{3}}, -\sqrt{\frac{2}{3}}\right)
    \end{equation*}
    Plugging them in we find the extremums are $\sqrt{6}$ and $-\sqrt{6}$.
\end{proof}


\begin{prob}[Marsden-Tromba, III. 16]
    Use Lagrange multipliers to find the distance from the point $(2, 0, -1)$ to the plane $3x - 2y + 8z + 1 = 0$.
\end{prob}
\begin{proof}
    Let $(x,y,z)$ be a point on the plane, then we would like the minimize the function 
    \begin{equation*}
        f(x,y,z)=(x-2)^2+y^2+(z+1)^2
    \end{equation*}
    with the constraint
    \begin{equation*}
        3x - 2y + 8z + 1 = 0
    \end{equation*}
    Writing $g(x,y,z)=3x - 2y + 8z + 1$, we do the exact same thing as we did above:
    \begin{equation*}
        \nabla f=\begin{bmatrix}
            2x-4&2y&2z+2
        \end{bmatrix}, \quad \nabla g=\begin{bmatrix}
            3&-2&8
        \end{bmatrix}
    \end{equation*}
    Then setting 
    \begin{align*}
        2x-4&=3\lambda\\
        2y&=-2\lambda\\
        2z+2&=8\lambda
    \end{align*}
    One can then replace $x,y,z$ with expressions in $\lambda$ in the constraint:
    \begin{equation*}
        3x-2y+8z+1=0
    \end{equation*}
    which gives 
    \begin{equation*}
        \lambda=2/77
    \end{equation*}
    Then plugging in $\lambda$ to solve for $x,y,z$, the $\sqrt{f(x,y,z)}$ is the final answer. (I am too lazy to do the computation).
\end{proof}







\end{document}